% EJERCICIO 9: corrige los 10 errores marcados con «% ERROR n».
% Compila con: latexmk -pdf ejercicio_apa7.tex
\documentclass[stu,12pt,a4paper,floatsintext]{apa7}
\usepackage[utf8]{inputenc}
\usepackage[T1]{fontenc}
\usepackage[spanish]{babel}   % ERROR 1: las tablas se rotulan «Cuadro»
\usepackage{booktabs}
\usepackage[style=apa,backend=biber,sortcites=true,sorting=nyt]{biblatex}
% ERROR 2: faltan csquotes y el mapeo de idioma de biblatex-apa
\addbibresource{ejercicio_apa7.bib}

\title{Acondicionamiento de señales con amplificadores operacionales}
% ERROR 3: falta el título corto (encabezado de página)
\authorsnames{Estudiante Ejemplo}
\authorsaffiliations{Universidad Nacional Tecnológica de Lima Sur}
\course{IE-205: Electrónica Analógica}
\professor{Docente Ejemplo}
\duedate{27 de octubre de 2026}
\abstract{Informe sobre amplificadores operacionales.}  % ERROR 4: el resumen debe tener 150-250 palabras
\keywords{amplificador operacional, ganancia, acondicionamiento de señales}

\AtBeginDocument{\renewcommand{\notelabel}{Nota}}
\begin{document}
\maketitle

\section{Introducción}
El amplificador operacional es un bloque fundamental del análisis de circuitos
analógicos \cite{sedra2020}.   % ERROR 5: falta el paréntesis APA
 Como señala \parencite{boylestad2023}   % ERROR 6: cita narrativa con el comando equivocado
, la ganancia
en lazo cerrado depende de la razón entre resistencias.

El formato de este informe sigue el manual de la \textcite{apa2020}.

\section{Resultados}
La Tabla 1 muestra   % ERROR 7: referencia escrita a mano
 la ganancia para tres valores de $R_f$.

\begin{table}[htbp]
  \caption{Ganancia del amplificador inversor}
  \label{tab:ganancia}
  \centering
  \begin{threeparttable}
  \begin{tabular}{|c|c|}   % ERROR 8: APA no usa líneas verticales ni \hline
    \hline
    $R_f$ (k$\Omega$) & Ganancia $-R_f/R_1$ \\
    \hline
    10 & $-1$ \\
    47 & $-4{,}7$ \\
    100 & $-10$ \\
    \hline
  \end{tabular}
  \begin{tablenotes}[para]
    \item \textit{Nota.} Con $R_1 = 10$~k$\Omega$.
  \end{tablenotes}
  \end{threeparttable}
\end{table}

\section{Uso de inteligencia artificial}
Se utilizó ChatGPT para   % ERROR 9: falta citar la herramienta
 depurar errores de compilación
del código \LaTeX. Las fuentes fueron verificadas por el autor.

\printbibliography[title={Referencias}]
\end{document}
